\section{Fazit}
Die Untersuchung zeigt, dass Large Language Models aus natürlichsprachlichen
Programmieraufgaben grundsätzlich geeignete Black-Box-Testfälle ableiten können. 
Auf Basis der Analyse kann die Forschungsfrage damit dahingehend beantwortet werden, 
dass ein LLM die Testfallerstellung wirksam unterstützen kann, die erzeugten Ergebnisse 
jedoch nicht ungeprüft übernommen werden sollten.

Bei der Nutzung der entwickelten Anwendung sollte eine möglichst eindeutige
Aufgabenbeschreibung mit allen Eingabegrenzen und dem erwarteten Rückgabeformat übergeben
werden. Anschließend sind die generierten Eingaben auf Einhaltung der Constraints und die
Sollwerte anhand der Spezifikation oder einer vertrauenswürdigen Referenzimplementierung zu
prüfen. Erst die validierten Testfälle sollten gegen die zu untersuchende Implementierung
ausgeführt werden. Die Anwendung ist somit als unterstützendes Werkzeug zur Erweiterung einer
Testsammlung zu verstehen und nicht als Ersatz für eine fachliche Kontrolle.

Für weiterführende Untersuchungen könnte das Experiment mit mehr Programmieraufgaben, mehreren
falschen Lösungen pro Aufgabe und wiederholten Generierungsdurchläufen durchgeführt werden.
Zusätzlich könnten unterschiedliche Modelle und Promptvarianten verglichen werden. Außerdem 
wäre es hilfreich das LLM zu jeder Testfallgenerierung eine kurze Argumentation asugeben zu lassen, 
um dessen Vorgehen nachvollziehbarer zu machen. Mittels dieser Erweiterungen könnten tiefere Einglicke 
in mögliche Zusammenhänge zwischen Validits- und Fehlererkunngsrate gewonnen und potentiell 
weitere Einflussfaktoren identifiziert werden.

Es ist anzunehmen das dabei die Unabhängigkeit zwischen Code- und Testfallgenerierung von besonderer Bedeutung bleibt. 
Wir das selbe Modell für beide Prozesse eingesetzt, könnten sich Fehlinterpretationen in Code und
Tests wiederholen und gegenseitig bestätigen. In zukünftigen Systemen oder Prozessen sollten Testfälle deshalb in
einem getrennten Kontext, ohne Zugriff auf konkrete Implementierungen und gegebenenfalls
durch ein unabhängiges Modell erzeugt und geprüft werden.